\documentclass[11pt]{article}
\usepackage[a4paper,margin=25mm]{geometry}
\usepackage{amsmath,amssymb,amsthm,hyperref}
\newtheorem{theorem}{Theorem}\newtheorem{lemma}{Lemma}
\newcommand{\E}{\mathbb E}
\title{An almost square-root rate for the interior endpoint profile}
\author{Complete theorem; independently reviewed}
\date{30 September 2026}
\begin{document}\maketitle

\begin{theorem}\label{main}
Let $p_1,\ldots,p_m$ be nondecreasing functions on a finite ordered
row set with positive probability weights $w$, and suppose
\[
 \E_w\prod_{j\in S}p_j=\frac1{|S|+1}\qquad(S\subseteq[m]),
 \qquad 0\le p_j\le1.
\]
Put $u_j=1-p_j$ and $s=\sum_j u_j$. For every fixed
$0<\delta<L<\infty$, uniformly over columns and rows with
$\delta\le s\le L$,
\[
 \max_j|m u_j-s|\le C_{\delta,L}\sqrt{\frac{\log m}{m}}
\]
for all sufficiently large $m$.
\end{theorem}

The constants are independent of the number of rows and of their
positive weights. The fixed lower cutoff $\delta>0$ is part of the
statement; no estimate at the endpoint $s=0$ is claimed here.

\paragraph{Inputs and their precise roles.}
This proof uses the independently reviewed notes
\begin{enumerate}
\item \path{../rational_designs/asymptotic/quantitative_uniform_endpoint_profile.tex}:
 on every fixed bounded $s$-interval, $\max_j|m u_j-s|\le m^{-\alpha}$
 for some fixed $\alpha>0$, after decreasing the exponent to absorb
 logarithms;
\item \path{marked_cancellation_identity.tex}: the exact marked
 identity, centered coefficient estimate, all-tilt subset
 concentration, and localization-preserving subset comparisons;
\item \path{../localized_laguerre_kernel_derivatives.tex}: the
 all-orders localized complex-derivative bounds for the positive
 Laguerre Poisson kernel;
\item \path{../rational_designs/asymptotic/global_laguerre_spectral_truncation.tex}:
 global marked and unmarked residuals after truncation at differential
 order $J=O(\log m)$, and their off-annulus tail bounds;
\item \path{../size_bounds/localized_profile_maximum_principle.tex}:
 the localized mass-absorption lemma. We spell out the final
 weighted maximum argument below to handle the different deleted
 column coordinates explicitly.
\end{enumerate}
No scalar transfer conjecture, rationality, or equal-weight assumption
enters the argument.

\paragraph{Coordinates and probability measures.}
Fix a marked column $i$, and set
\[
 M=m-1,\quad \ell=\lfloor M/4\rfloor,\quad N=M-\ell,
 \quad z_j=\ell u_j.
\]
For a uniformly selected $\ell$-subset $B$ of the other $M$ indices,
let $R$ be its complement and put
\[
 A_B=\prod_{j\in B}p_j,\quad
 v_B=\frac1N\sum_{j\in R}z_j,\quad
 x_i=\frac1M\sum_{j\ne i}z_j,\quad h_i=z_i-x_i.
\]
Write $\overline A_i=\E_BA_B$ and define the row measure
\[
 d\mu_i=(\ell+1)\overline A_i\,dw.
\]
It has total mass one by the mixed-moment identities.
The scalar reference probability measure is
\[
 d\beta_\ell(x)=\frac{\ell+1}{\ell}
 (1-x/\ell)^\ell\boldsymbol1_{[0,\ell]}(x)\,dx.
\]
Its density is bounded above and below on every fixed positive
compact interval, uniformly for large $m$.

For a fixed large exponent $A$, use the choices in the global
spectral-truncation theorem:
\[
 \tau=C_A\frac{\log m}{m},\qquad b=\sqrt\tau,
 \qquad J=\lceil B_A\log m\rceil.
\]
All constants are fixed before taking $m$ large. Write $K_\tau(x,a)$
for the positive Laguerre Poisson kernel. On positive compact
intervals it has Gaussian upper bounds at scale $b$ and a lower
bound $c/b$ whenever $|x-a|\le b$.

\begin{lemma}[The two actual localized comparisons]\label{comparisons}
Fix a compact positive interval of test centers. There is a larger
fixed compact interval $W\subset(0,\infty)$ of row coordinates such
that, uniformly over marked columns and test centers, the following
bounds hold. Let
$H(q)=\max_j z_j(q)-\min_j z_j(q)$, including the marked column.
Then
\begin{align}
 |\mu_i(K_\tau(x_i,a))-\beta_\ell(K_\tau(\cdot,a))|
 &\le \frac C{m\tau}\int_{x_i\in W}
                H^2K_{c\tau}(x_i,a)\,d\mu_i+m^{-A},\label{unmarked}\\
 |\mu_i(h_iK_\tau(x_i,a))|
 &\le \frac C m\int_{x_i\in W}
 H^2(1+b^{-1}+Hb^{-2})K_{c\tau}(x_i,a)\,d\mu_i+m^{-A}.
 \label{marked}
\end{align}
Here $c$ is a fixed enlargement factor. On $W$ one has
$H\le R_m\le m^{-\alpha}$ for some fixed $\alpha>0$.
Kernel tails outside $W$, with any fixed polynomially bounded
prefactors used in these formulas, are negligible.
\end{lemma}
\begin{proof}
Choose $W$ from the off-annulus theorem in the global truncation
note, enlarging it if necessary. Its row coordinate there is
$s_{-i}=\sum_{j\ne i}u_j=(M/\ell)x_i$; the conversion changes only
fixed endpoints. On $W$, the finite balance bound
$z_i\le C(s_{-i}+1)$ makes the full $s$ bounded. The prior profile
theorem therefore gives $H\le m^{-\alpha}$ after weakening a power.
All $z_j$ then lie in a fixed positive compact interval, since their
unmarked mean is $x_i$, and $p_j\ge1/2$ for sufficiently large $m$.

The global unmarked differential residual compares the reference
integral to
\[
 (\ell+1)\E_w\E_B A_B
 \sum_{k=0}^J \frac{K_\tau^{(k)}(v_B,a)}{k!}a_k,
 \qquad
 a_k=\frac{e_k((z_j-v_B)_{j\in R})}{\binom Nk}.
\]
The off-annulus theorem removes rows outside $W$ with error
$m^{-A}$, even with absolute values on each summand. On $W$, the
centered local derivative estimate bounds the terms $k\ge2$ by
$CH^2/(m\tau)$ times a broader kernel. Average under the
success-weighted subset law, then use its all-tilt concentration to
keep that envelope at $x_i$. The localized subset Taylor comparison
replaces $K_\tau(v_B,a)$ by $K_\tau(x_i,a)$ with the same bound.
This proves \eqref{unmarked}.

For the marked identity the global residual is
\[
 (\ell+1)\E_w\E_B A_B\left[(z_i-v_B)K_\tau(v_B,a)
 -\sum_{k=2}^J\left\{
 \frac{(v_B-z_i)K_\tau^{(k)}(v_B,a)}{k!}
 +\frac{K_\tau^{(k-1)}(v_B,a)}{(k-1)!}\right\}a_k\right].
\]
It is at most $m^{-A}$ in absolute value. On $W$, the finite
differential sum is bounded by
$C[H^2/(mb)+H^3/(mb^2)]$ times a broader kernel, by the direct
marked kernel lemma. The localized subset comparison replaces
$(z_i-v_B)K_\tau(v_B,a)$ by $h_iK_\tau(x_i,a)$ with error
$CH^2(1+b^{-1}+Hb^{-2})/m$ times such a kernel. This proves
\eqref{marked}.

The pure scalar kernels outside $W$ have exponentially small
Gaussian tails because their centers lie in the smaller fixed
interval. Every row measure $\mu_i$ has mass one, and the row
coordinates and the auxiliary factors are at most polynomial in
$m$. Thus these scalar tails are negligible too. Increasing the
initial exponent $A$ absorbs the fixed number of residuals.
\end{proof}

\paragraph{Uniform local mass.}
On $W$, $H\le R_m$, so \eqref{unmarked} is bounded by
\[
 \epsilon_m\mu_i(K_{c\tau}(x_i,a))+m^{-A},
 \qquad \epsilon_m=\frac{CR_m^2}{m\tau}\longrightarrow0.
\]
Apply the mass-absorption lemma on an enlarged fixed annulus.
Consequently, on any prescribed smaller annulus $I$, uniformly in
$i$ and the center $a$,
\begin{equation}\label{mass}
 \mu_i(|x_i-a|\le b)\le Cb,\qquad
 c\le\mu_i(K_\tau(x_i,a))\le C.
\end{equation}
Broader kernel masses are also bounded. Gaussian integration over
length-$b$ intervals gives, for $r\ge b$ within that annulus,
\begin{equation}\label{tail}
 \int_{|x_i-a|\ge r}K_\tau(x_i,a)\,d\mu_i
       \le C e^{-cr^2/b^2}+O(m^{-A}).
\end{equation}
The outer annulus for this step is chosen strictly larger than $I$,
so the constants remain uniform up to the boundaries of $I$.

\paragraph{A common coordinate for all marked columns.}
Put
\[
 y=\frac\ell m s,\qquad h_i^*=z_i-y.
\]
The exact identities
\begin{equation}\label{conversion}
 x_i=y-\frac{h_i^*}{M},\qquad
 h_i=\frac mM h_i^*,\qquad
 H=\max_j h_j^*-\min_j h_j^*
\end{equation}
will avoid comparing maxima in different deleted coordinates.
Choose a fixed positive annulus $I=[a_-,a_+]$, and a strictly
smaller $I_0$, so that all target rows $\delta\le s\le L$ have
$y\in I_0$ for large $m$. On a fixed enlargement of $I$ the prior
profile estimate gives $|h_i^*|\le R_m\le m^{-\alpha}$.
Let $\phi(y)=\min(y-a_-,a_+-y)$ on $I$, and set
\[
 D=\max_{i,\,\text{rows with }y\in I}\phi(y)|h_i^*|.
\]
If the row set here is empty, the target assertion is vacuous.

Suppose $D>C_*b$. Select a maximizing row and column $i$, write
$y_0$ for its common coordinate, $\Delta=|h_i^*(y_0)|$, and
$d_0=\phi(y_0)$. Then
\[
 d_0=D/\Delta\ge C_*b/R_m,\qquad d_0/b\longrightarrow\infty
\]
at a polynomial rate. Also $\Delta\ge c_I C_*b$.
For every row with $|y-y_0|\le d_0/2$, weighted maximality implies
\begin{equation}\label{localrange}
 |h_j^*(y)|\le C\Delta\quad\hbox{for all }j,
 \qquad H(y)\le C\Delta.
\end{equation}
Any row whose $x_i$ lies in the fixed test annulus has bounded
$s_{-i}$, and finite balance then makes its full $s$ bounded. Thus it
lies in a fixed $y$-enlargement where the prior profile estimate
applies, even if it was not initially described by its $y$ coordinate.
Moreover, $|x_i-y|\le R_m/M=o(b)$ on the entire enlarged annulus,
by \eqref{conversion}. Thus \eqref{localrange} holds on a fixed
smaller $x_i$-neighborhood of the selected deleted coordinate
$x_0=y_0-h_i^*(y_0)/M$. The cutoff weights there are comparable to
$d_0$. Rows not in this neighborhood are separated from $x_0$ by
a fixed multiple of $d_0$, apart from negligible $o(b)$ shifts.

\paragraph{A one-sided positive test.}
Choose fixed $B_0$ large enough for the Gaussian tails in
\eqref{tail}, and then set $A_0=2B_0$. If the selected deviation is
positive, center the kernel at $a=x_0+A_0b$; if it is negative,
take $a=x_0-A_0b$. Both centers remain inside $I$.
Consider the positive case. Whenever $|x_i-a|\le B_0b$, one has
$x_i>x_0$ strictly. The monotone ordering of the columns then gives
$z_i\ge z_i(x_0)$, even if there are other flat rows. Therefore
\[
 h_i(x_i)\ge h_i(x_0)-(x_i-x_0)
 \ge\frac mM\Delta-(A_0+B_0)b\ge\Delta/2,
\]
provided $C_*$ was chosen large enough. For a negative selected
deviation the same argument to the left gives $h_i\le-\Delta/2$.

The kernel mass on this favorable window is bounded below by a
positive constant using \eqref{mass}--\eqref{tail}, after choosing
$B_0$ sufficiently large. In the surrounding $d_0/4$ neighborhood,
\eqref{localrange} bounds $|h_i|$ by $C\Delta$, so its opposite-sign
contribution outside the favorable window is as small a fixed
multiple of $\Delta$ as desired. Outside that neighborhood the
Gaussian factor is $\exp[-c(d_0/b)^2]$. Since $d_0/b$ grows
polynomially, every polynomial row factor is absorbed. It follows
that
\begin{equation}\label{positive}
 |\mu_i(h_iK_\tau(x_i,a))|\ge c\Delta.
\end{equation}

On the other hand, \eqref{marked}, the local range bound, bounded
broader kernel mass, and the same negligible far tails give
\[
 |\mu_i(h_iK_\tau(x_i,a))|
 \le\frac{C\Delta^2}{m}(1+b^{-1}+\Delta b^{-2})+m^{-A}
 =o(\Delta).
\]
Indeed $mb^2=C_A\log m$, $mb\to\infty$, and
$\Delta\le R_m\to0$, whereas $\Delta\ge c_I C_*b$.
This contradicts \eqref{positive}. Hence $D\le C_*b$.
On $I_0$ the cutoff is bounded below, so $|h_i^*|\le Cb$.
Finally $m u_i-s=(m/\ell)h_i^*$ and $m/\ell$ is bounded.
This proves Theorem~\ref{main}.
\end{document}
